% a mashup of hipstercv, friggeri and twenty cv
% https://www.latextemplates.com/template/twenty-seconds-resumecv
% https://www.latextemplates.com/template/friggeri-resume-cv

\documentclass[lighthipster]{simplehipstercv}
% available options are: darkhipster, lighthipster, pastel, allblack, grey, verylight, withoutsidebar
% withoutsidebar
\usepackage[utf8]{inputenc}
\usepackage[default]{raleway}
\usepackage[margin=1cm, a4paper]{geometry}


%------------------------------------------------------------------ Variablen

\newlength{\rightcolwidth}
\newlength{\leftcolwidth}
\setlength{\leftcolwidth}{0.23\textwidth}
\setlength{\rightcolwidth}{0.75\textwidth}

%------------------------------------------------------------------
\title{New Simple CV}
\author{\LaTeX{} Ninja}
\date{June 2019}

\pagestyle{empty}
\begin{document}


\thispagestyle{empty}
%-------------------------------------------------------------

\section*{Start}

\simpleheader{headercolour}{Leandro}{Olguin}{Estudiante de Ingeniería Industrial}{white}


%------------------------------------------------

% this has to be here so the paracols starts..
\subsection*{}
\vspace{4em}

\setlength{\columnsep}{1.5cm}
\columnratio{0.23}[0.75]
\begin{paracol}{2}
\hbadness5000
%\backgroundcolor{c[1]}[rgb]{1,1,0.8} % cream yellow for column-1 %\backgroundcolor{g}[rgb]{0.8,1,1} % \backgroundcolor{l}[rgb]{0,0,0.7} % dark blue for left margin

\paracolbackgroundoptions

% 0.9,0.9,0.9 -- 0.8,0.8,0.8


\footnotesize
{\setasidefontcolour
\flushright
\begin{center}
    \roundpic{fotop.jpg}
\end{center}

\bg{cvgreen}{white}{Sobre mí}\\[0.5em]

{\footnotesize
"Estudiante de Ingeniería Industrial con una sólida base técnica en informática y gran capacidad para la resolución de problemas complejos mediante el análisis matemático y estadístico. Destaco mi gran responsabilidad y habilidad para el aprendizaje de nuevas tecnologías.}
\bigskip

\bg{cvgreen}{white}{Personal} \\[0.5em]
Leandro Olguin

22 años

Mendoza, Argentina


\bigskip

\bg{cvgreen}{white}{Áreas de especializaciónn} \\[0.5em]

~•~ Física \\
~•~ Geometría Analítica\\
~•~ Química\\ 
~•~ Electrotecnia \\

\bigskip



\bigskip

\bg{cvgreen}{white}{Intereses}\\[0.5em]
~•~ Optimización de Procesos \\ 
~•~ Análisis de Datos Técnicos \\
~•~Automatización y Control\\
~•~Tecnologías Emergentes\\
\bigskip


\vspace{4em}

\infobubble{\faAt}{cvgreen}{white}{\parbox[t]{6.6em}{leannnolguin\\@gmail.com}}
\infobubble{\faGithub}{cvgreen}{white}{Leandro-olguin}

\phantom{turn the page}

\phantom{turn the page}
}
%-----------------------------------------------------------
\switchcolumn

\small
\section*{Formación académica}

\begin{tabular}{r| p{0.4\textwidth} c}
    \cvevent{2022--Actualidad}{Ingenieria industrial}{Estudiante}{Universidad Nacional de Cuyo \color{cvred}}{Formacion en gestión de operaciones, optimización de recursos y análisis de sistemas productivos. Conocimientos avanzados en cálculo de varias variables y ecuaciones diferenciales. Desarrollo de algoritmos en \textbf{MATLAB} y \textbf{Octave} para la resolución de sistemas complejos mediante métodos numéricos (Jacobi, Runge-Kutta, Euler).}{unc.jpg} \\
    \cvevent{2015--2020}{Bachiller en informática}{Colegio Martin Zapata}{Mendoza \color{cvred}}{Desarrollo de habilidades en lógica de programación, gestión de bases de datos y soporte técnico. Orientación técnica que facilita el manejo de software especializado de ingeniería y herramientas de automatización de oficina.}{mz.jpg}
\end{tabular}
\vspace{3em}

\begin{minipage}[t]{0.35\textwidth}
\section*{Proyectos y experiencia}
\begin{tabular}{r p{0.6\textwidth} c}
    \cvdegree{2022}{Introducción a la ingenieria}{Informe sobre empresa de ensayos no destructivos}{UNcuyo \color{headerblue}}{}{unc.jpg} \\
    \cvdegree{2025}{Electrotecnia}{Diseño y Simulación de Circuitos}{UNCuyo \color{headerblue}}{}{unc.jpg} \\
    \cvdegree{2025}{Termodinámica}{Análisis de Ciclos de Potencia}{UNCuyo \color{headerblue}}{}{unc.jpg}
\end{tabular}
\end{minipage}\hfill
\begin{minipage}[t]{0.3\textwidth}
\section*{Habilidades técnicas}
\begin{tabular}{r @{\hspace{0.5em}}l}
     \bg{skilllabelcolour}{iconcolour}{html} &  \barrule{0.4}{0.5em}{cvpurple}\\
     \bg{skilllabelcolour}{iconcolour}{\LaTeX} & \barrule{0.55}{0.5em}{cvgreen} \\
     \bg{skilllabelcolour}{iconcolour}{Octave} & \barrule{0.5}{0.5em}{cvpurple} \\
     \bg{skilllabelcolour}{iconcolour}{Python} & \barrule{0.25}{0.5em}{cvpurple} \\
     \bg{skilllabelcolour}{iconcolour}{R} & \barrule{0.03}{0.5em}{cvpurple} \\
     
\end{tabular}
\end{minipage}

\section*{Curriculum}
\begin{tabular}{r| p{0.5\textwidth} c}
    \cvevent{2025}{Electrotecnia}{Análisis de Sistemas de Energía}{UNCuyo \color{cvred}}{Estudio de eficiencia en redes eléctricas y máquinas térmicas, con enfoque en la reducción de pérdidas energéticas en sistemas industriales.}{unc.jpg} \\
    \cvevent{2025}{Termodinámica}{Simulación de Ciclos Térmicos}{UNCuyo \color{cvred}}{Cálculo de balances de masa y energía en sistemas abiertos y cerrados. Análisis detallado de ciclos de potencia (Rankine/Otto) y optimización de la eficiencia térmica mediante el uso de tablas de vapor y diagramas de fases.}{unc.jpg} \\
\end{tabular}
\vspace{3em}

\begin{minipage}[t]{0.3\textwidth}
\section*{Logros académicos}
\begin{tabular}{>{\footnotesize\bfseries}r >{\footnotesize}p{0.55\textwidth}}
    2020 & Egresado con honores - Bachiller en Informática \\
    2021 & Especialización en lenguajes de programación \\
    2026 & Presentación académica (UNCuyo) 
\end{tabular}
\bigskip

\section*{Languajes}
\begin{tabular}{l | ll}
\textbf{Español} & Nativo & \pictofraction{\faCircle}{cvgreen}{4}{black!30}{1}{\tiny} \\
\textbf{Inglés} & Técnico & \pictofraction{\faCircle}{cvgreen}{3}{black!30}{2}{\tiny}\\
\textbf{Italiano} & Básico & \pictofraction{\faCircle}{cvgreen}{1}{black!30}{4}{\tiny} \\
\end{tabular}
\bigskip

\end{minipage}\hfill
\begin{minipage}[t]{0.3\textwidth}
\section*{Proyectos}
\begin{tabular}{>{\footnotesize\bfseries}r >{\footnotesize}p{0.7\textwidth}}
    2019 & \emph{Página web para la publicación de información útil en el cursado y comunicados}\\
    2026 & ``Bitacora de aprendizaje'', desarrollo de un perfil profesional detallado y registro de procesos técnicos en ingeniería relacionados con el trato hacia las personas.
\end{tabular}
\bigskip

\end{minipage}






\vfill{} % Whitespace before final footer

%----------------------------------------------------------------------------------------
%	FINAL FOOTER
%----------------------------------------------------------------------------------------
\setlength{\parindent}{0pt}
\begin{minipage}[t]{\rightcolwidth}
\begin{center}\fontfamily{\sfdefault}\selectfont \color{black!70}
{\small Leandro Olguin \icon{\faEnvelopeO}{cvgreen}{} Mendoza \icon{\faMapMarker}{cvgreen}{} Argentina \icon{\faPhone}{cvgreen}{} 2615458037 \newline\icon{\faAt}{cvgreen}{} \protect\url{leannnolguin@gmail.com}
}
\end{center}
\end{minipage}

\end{paracol}

\end{document}
